\section{Inference and Trajectory Generation}

Once the model is trained, it can be used as a high-fidelity function approximator to generate entire trajectories on demand.

\subsection{Generating a Trajectory}
To predict the evolution of a system from a given initial condition ($IC$), the following procedure is used:
\begin{enumerate}
    \item \textbf{Define Evaluation Times:} Create a fine-grained array of time points, $t_{eval}$, over the desired simulation interval (e.g., 10,000 points from $t=0$ to $t_{max}$). Normalize these time points to the model's input range of $[-1, 1]$.
    \item \textbf{Predict Positions:} For each normalized time point in $t_{eval}$, feed it concatenated with the constant $IC$ vector into the trained network. This will produce a continuous trajectory of the system's positions, $r_{\theta}(t_{eval})$.
    \item \textbf{Derive Kinematics (Optional):} If velocities or accelerations are needed for analysis, they can be derived at any evaluation point by using automatic differentiation on the network with respect to the time input.
\end{enumerate}





\subsection{Output Layer and Kinematic Consistency via Automatic Differentiation}


The network architecture is simplified and physical consistency is hard-coded by predicting only positions and deriving all other kinematic quantities via automatic differentiation (AD). This is a core design principle of PINNs, supported by established frameworks like DeepXDE.$^{14}$

The final layer of the SQR-SkipResNet-SIREN network is a standard linear layer without any activation function. This is crucial as it allows the network's outputs—the predicted positions—to take on any real value required by the physical system. The output is a 9-dimensional vector representing the concatenated positions of the three bodies at a given time $\mathrm{t}$:
\[
\mathrm{r}_\theta(\mathrm{t}) = [\mathrm{r}_1(\mathrm{t}), \mathrm{r}_2(\mathrm{t}), \mathrm{r}_3(\mathrm{t})]
\]

Crucially, the network does not predict velocities or accelerations directly. Instead, these kinematic quantities are derived from the position outputs using AD with respect to the time input, $t$. The velocity vector is computed as the first-order derivative of the position vector, and the acceleration vector is computed as the second-order derivative$^1$:
\[
\mathrm{v}_\theta(\mathrm{t}) = \frac{\mathrm{d} \mathrm{r}_\theta(\mathrm{t})}{\mathrm{d} \mathrm{t}} \quad \text{and} \quad \mathrm{a}_\theta(\mathrm{t}) = \frac{\mathrm{d}^2 \mathrm{r}_\theta(\mathrm{t})}{\mathrm{d} \mathrm{t}^2}
\]

By computing derivatives via AD, the fundamental kinematic relationships between position, velocity, and acceleration $(\mathrm{v} = \mathrm{r}', \mathrm{a} = \mathrm{r}'')$ are satisfied by construction to machine precision. This hard-codes a significant portion of the problem's physics into the model, simplifies the network architecture (as no extra heads are needed for velocity or acceleration), and makes the overall learning task more tractable.





\subsection{Output Frame and File Format}
For consistency and physical correctness, the output of the model must be clearly defined.
\begin{itemize}
    \item \textbf{Reference Frame:} The predicted positions, $r_{\theta}(t)$, are in the \textbf{Center of Mass (COM) frame}. This is consistent with the use of translation-invariant quantities in the input vector.
    \item \textbf{Output File Format:} The final generated trajectory should be saved in a structured format, such as a CSV file, for easy analysis and visualization. Each row should represent a single time step.
    \item \textbf{Output Definition:} The output at each network forward pass is strictly the positions of the three bodies at a single normalized input time $t'$;
    velocities and accelerations are derived via AD.
\end{itemize}

\textbf{Example CSV Output Format:}
\begin{verbatim}
time,x1,y1,z1,x2,y2,z2,x3,y3,z3
0.00,-0.98,0.20,0.00,0.49,-0.10,0.00,0.49,-0.10,0.00
0.01,-0.97,0.21,0.00,0.48,-0.11,0.00,0.49,-0.10,0.00
...
\end{verbatim}
